\chapter{Testing and Evaluation}\label{ch:testing}

All results in this chapter come from commands executed on the reference laptop (\autoref{tab:env}) between 8 and 9
October 2026, or from the project's continuous-integration runs. Evaluation numbers are taken from run records
committed under \code{eval/runs/}, which store the git commit, hardware, power source, dataset hash and full
configuration. Where a measurement has a caveat, the caveat is stated with it.

\section{Test Strategy}

Tests are organised in layers, selected with pytest markers so that each layer runs where its dependencies exist.
Tests that replace a component with a fake say so in their module documentation.

\begin{table}[htbp]
  \centering
  \caption{Test layers and results in the final acceptance run (9 October 2026)}\label{tab:testlayers}
  {\small
  \begin{tabularx}{\textwidth}{@{}lLlr@{}}
    \toprule
    \textbf{Layer} & \textbf{What is real / what is faked} & \textbf{Runs in CI} & \textbf{Passed} \\
    \midrule
    Backend unit & Pure code; fake embedder, scripted fake model, fake arXiv & Yes & 140 \\
    Integration and model & Real PostgreSQL + pgvector and real MiniLM; fake model & Yes & 83 \\
    Real local model & Ollama \code{qwen2.5:3b} & No & 3 \\
    Live arXiv & arXiv API & No & 1 \\
    Stack end-to-end & Running Docker stack through nginx, all real components & No & 7 \\
    Frontend & Vitest and Testing Library with a mocked HTTP client & Yes & 24 \\
    \bottomrule
  \end{tabularx}
  }
\end{table}

Static checks (Ruff lint and format, mypy in strict mode over 78 files, oxlint, the TypeScript compiler) passed, and
the production frontend build succeeded. The live arXiv test failed once while arXiv was returning HTTP 429 to the
machine and passed when re-run; it is reported as such. No test was skipped in the integration layer: in CI a
missing test database makes the run fail rather than skip.

\subsection{Unit, API and Integration Tests}

Unit tests cover parsing of arXiv identifiers and feeds, query sanitisation, PDF extraction and page mapping,
deterministic chunking, configuration validation, prompt construction, the output schema, citation checking
(invalid, duplicate and label-only citations), error mapping of the Ollama adapter, the evaluation metrics, search
mode routing and fusion, and the error envelope. API tests use FastAPI's test client to check input bounds,
pagination limits, CORS, request identifiers and readiness. Integration tests run the real pipeline against
PostgreSQL: arXiv import (with a fake arXiv source), idempotent re-import, upload validation, duplicate detection,
search in all four modes, question answering with a scripted model (including fabricated labels, abstention and
generation failures), analyses, corpus management, the evaluation harness, and the provenance checker.

\subsection{Database and Migration Tests}

Integration tests apply all migrations to an empty database, run \code{alembic check} to detect drift between the
ORM models and the migrations, perform a downgrade and upgrade round trip, and verify constraints such as cascading
deletes and the vector dimension. A reproducibility check on a fresh clone migrated a new database from revision
0001 to 0006 on first start.

\subsection{End-to-End Tests}

The end-to-end suite runs against the Docker stack: it checks readiness, uploads a synthetic paper with a unique
marker word, waits for ingestion, finds the paper with all four search modes, asks a question and checks that a
citation is reported valid exactly when its label is among the evidence, requests a summary, and deletes the paper.
During final acceptance the same journeys were also exercised manually with real papers
(\autoref{sec:acceptance-journeys}).

\section{Evaluation Data Set}

\begin{itemize}
  \item \textbf{Corpus:} 20 natural-language-processing papers pinned to exact arXiv versions
        (\code{eval/corpus.json}); 399 pages and 1{,}654 chunks.
  \item \textbf{Questions:} \code{eval/qa\_v1.json}, 50 items: 40 answerable with 110 verbatim evidence quotes, and 10
        unanswerable, several of them adversarial (they name a corpus paper and ask for a fact found only in a
        neighbouring paper).
  \item \textbf{Provenance of labels:} all items were written by the AI assistant used by the team (ADR-0007) by
        keyword search of the papers' text, not with the retriever. \textbf{No item has been reviewed by a person.}
        Each quote is verified to occur in the corpus before every run.
  \item \textbf{Relevance unit:} a retrieved chunk is relevant if it belongs to the labelled paper and contains one
        of the item's evidence quotes.
\end{itemize}

\section{Metrics}

For $Q$ answerable questions, let $r_q$ be the rank of the first relevant chunk for question $q$ (undefined if none
is in the top 10). Then
\begin{equation}
  \recall{k} = \frac{1}{Q}\sum_{q=1}^{Q} \mathbf{1}[r_q \leq k], \qquad
  \mathrm{MRR} = \frac{1}{Q}\sum_{q=1}^{Q} \frac{1}{r_q},
\end{equation}
with $1/r_q = 0$ when no relevant chunk is retrieved. Recall@$k$ is thus a hit rate, as in the proposal. Both are
implemented in \code{backend/app/evaluation/metrics.py} and tested on hand-computed cases. Citation validity is the
number of valid citations divided by all citations in answers returned as answered. Abstention recall is the
fraction of unanswerable questions that receive ``insufficient evidence''; abstention precision is the fraction of
``insufficient evidence'' outcomes that occur on unanswerable questions. Latency percentiles use the nearest-rank
method.

\section{Retrieval Results}

\autoref{tab:retrieval} compares the four search modes. Configuration: top 10, RRF $k=60$, candidate pool
$\max(4k,20)$ per list, AC power, CPU inference.

\begin{table}[htbp]
  \centering
  \caption{Retrieval quality and search latency by mode (40 answerable questions)}\label{tab:retrieval}
  {\small
  \begin{tabular}{@{}lrrrrr@{}}
    \toprule
    \textbf{Mode} & \textbf{R@1} & \textbf{R@5} & \textbf{R@10} & \textbf{MRR} & \textbf{p50 / p95 (ms)} \\
    \midrule
    dense (baseline) & 0.200 & 0.700 & 0.775 & 0.416 & 42 / 49 \\
    bm25 & 0.300 & 0.525 & 0.650 & 0.393 & 56 / 91 \\
    hybrid (RRF) & 0.375 & 0.725 & 0.825 & 0.527 & 102 / 137 \\
    hybrid\_rerank (default) & \textbf{0.575} & \textbf{0.875} & \textbf{0.925} & \textbf{0.701} & 1{,}734 / 1{,}955 \\
    \bottomrule
  \end{tabular}
  }

  \smallskip
  {\footnotesize Run records: dense and hybrid\_rerank \code{20261008T214635Z} and \code{20261008T214832Z}
  (commit \code{03fb705}); bm25 and hybrid \code{20261008T203750Z} and \code{20261008T203825Z} (commit \code{192ac2b}).}
\end{table}

Only \code{hybrid\_rerank} meets the target Recall@5 $\geq 0.80$ (NFR-01). Three caveats apply: the mode was chosen
on the same 50 questions it is measured on; the labels are unreviewed; and with 40 questions the 95\% Wilson interval
for 35/40 is roughly 0.74 to 0.95, which includes values below the target. Every miss in every mode retrieved the
correct paper but ranked the answering passage too low. Metrics were identical across four commits, which confirms
that retrieval is deterministic.

\section{Retrieval Latency Profile}

\autoref{tab:profile} breaks retrieval into stages for the question-answering configuration (top 6, 24 candidates),
measured with the project's \code{profile-retrieval} command over the 50 questions.

\begin{table}[htbp]
  \centering
  \caption{Retrieval stage latency at top 6 (ms, p50 / p95)}\label{tab:profile}
  {\small
  \begin{tabular}{@{}lr@{}}
    \toprule
    \textbf{Stage} & \textbf{Latency} \\
    \midrule
    Query embedding (MiniLM, CPU) & 17 / 30 \\
    Dense SQL (exact scan) & 29 / 36 \\
    Lexical SQL & 36 / 66 \\
    Hybrid call (dense + lexical + RRF) & 61 / 85 \\
    Cross-encoder re-ranking (24 pairs, CPU) & 895 / 985 \\
    \bottomrule
  \end{tabular}
  }
\end{table}

The cross-encoder accounts for about 92\% of retrieval time. The same profile found 47 pairs of neighbouring
(overlapping) chunks among the 300 top-6 slots and no exact duplicates. At this corpus size the PostgreSQL planner
uses an exact sequential scan (12\,ms) instead of the HNSW index, so results are exact.

\section{Question-Answering Results}

Two complete runs over all 50 questions used the default retrieval, \code{qwen2.5:3b} at temperature 0, top 6 and the
0.30 threshold (\autoref{tab:qa}). The first run also scored each claim with a local judge.

\begin{table}[htbp]
  \centering
  \caption{Grounded question answering over the evaluation set}\label{tab:qa}
  {\small
  \begin{tabularx}{\textwidth}{@{}Lrr@{}}
    \toprule
    \textbf{Measure} & \textbf{Run 1 (\code{204611Z})} & \textbf{Run 2 (\code{205117Z})} \\
    \midrule
    Answer rate on answerable questions & 0.875 & 0.950 \\
    Answerable questions with a valid citation to a relevant passage & 0.725 & 0.750 \\
    Citation validity (valid / all) & 52 / 52 & 57 / 57 \\
    Unanswerable questions declined (abstention recall) & 9 / 10 & 9 / 10 \\
    Abstention precision & 9 / 14 & 9 / 11 \\
    False answers on unanswerable questions & 1 / 10 & 1 / 10 \\
    Warm latency p50 / p95 (s) & 3.10 / 4.50 & 3.06 / 4.43 \\
    Cold start (three samples, s) & 8.4--12.0 & 8.1--10.2 \\
    \bottomrule
  \end{tabularx}
  }
\end{table}

The false answer is the same in both runs: asked how much energy training DistilBERT consumed, the model answered
with a figure from the LLaMA paper, citing a real passage about the wrong model. Run-to-run differences in the
answer rate show that generation is not bit-reproducible even at temperature 0; results are therefore reported per
run. Stored answers from run 2 give a median of 1.16\,s for retrieval and 1.95\,s for generation per question.

\section{Citation Integrity}

Citation integrity was checked deterministically, independently of the model:
\begin{itemize}
  \item across all 418 answers stored in the evaluation database, every one of 422 valid citations resolves to an
        evidence row of the same answer, and no invalid citation is linked to evidence;
  \item all 2{,}364 evidence snapshots equal the text, pages and paper of the chunk they came from;
  \item the provenance checker re-extracted all 20 PDFs and confirmed for all 1{,}654 chunks that the stored text
        equals the document text at the recorded span, that the recorded pages are correct, and that the chunk's
        first and last words occur on those pages of the PDF (0 problems; 433 chunks span two or more pages).
\end{itemize}
Unit and integration tests cover non-existent labels, duplicate labels, label-only claims, malformed model output
and deleted sources (evidence snapshots are kept with an empty chunk reference). Because the model can only emit
labels, a fabricated page number or chunk identifier cannot enter the system.

\section{Groundedness}

A valid citation shows that a claim refers to a passage the model was given, not that the passage supports the
claim. Two complementary checks were made:
\begin{itemize}
  \item \textbf{Local judge} (\code{judge-v1}, the same \code{qwen2.5:3b} model): of 37 claims in run 1, 31 (83.8\%)
        were rated supported, 3 partially supported and 3 not supported. The judge is not independent of the
        answering model.
  \item \textbf{Spot check by the AI assistant}, not by a person: 14 claims (the 6 the judge flagged and 8 random
        claims it rated supported) were read against the full cited text. Ten were supported and answered the
        question; one was supported but did not answer it; one was a wrong-paper answer (a Transformer question
        answered with LLaMA's optimiser settings); two were not stated in the passage they cited. The judge agreed
        with this reading on 11 of 14 claims.
\end{itemize}
\textbf{No human spot check has been performed}; it is required by the project's evaluation rules and remains open.

\section{Abstention, Failure Handling and Prompt Injection}

The system abstains at three points: when no passage reaches the threshold (no model call), when the model reports
insufficient evidence, and when no claim has a valid citation. In the evaluation runs all abstentions on
unanswerable questions came from the model. Failure handling was tested live: with the database stopped, readiness
reports the cause and the API returns 503 with an actionable message, recovering when the database restarts; with
Ollama unreachable, question answering returns 503 while search keeps working; invalid input returns 422 with the
offending field. A real-model test placed an instruction inside a retrieved passage; the model did not follow it
(\code{test\_instructions\_inside\_passages\_are\_not\_followed}).

\section{Final Acceptance Journeys}\label{sec:acceptance-journeys}

On 9 October 2026 the rebuilt stack was exercised with real components: discovering papers through arXiv search;
importing \code{2004.04906v3} (and confirming that a second import created no duplicate); uploading a real five-page
paper; re-checking provenance for all four stored documents (210 chunks, 0 problems); confirming 384-dimensional
embeddings for all chunks; searching; asking ``How much smaller and faster than BERT is DistilBERT?'' (answered
``40\% smaller and 60\% faster'' with two citations whose passages contain both figures, one spanning pages 4--5);
asking an unsupported question (declined before any model call); running summary, extraction, synthesis and
comparison; filtering and deleting papers; and opening citations in the user interface. One defect found during
this run (arXiv title search with stop words returned nothing) was fixed and covered by a regression test.

\section{Continuous Integration}

GitHub Actions runs six jobs on every push: backend lint, types and unit tests; integration tests against a pgvector
service container with the real embedding model; frontend checks and build; Docker image build; a retrieval
regression gate that rebuilds the evaluation corpus from arXiv and fails if Recall@5 drops below 0.85 for
\code{hybrid\_rerank} or 0.675 for \code{dense} (each one question below its baseline); and dependency scanning with
\code{pip-audit} and \code{npm audit}. One run on the main branch failed because arXiv rate-limited the runner;
the corpus builder now retries failed downloads with back-off. CI does not use Ollama, so real-model and end-to-end
tests are run manually.
